% =====================================================================
%  Synthèse — fiche en une page
% =====================================================================
\chapter*{Synthèse en une page}
\addcontentsline{toc}{chapter}{Synthèse en une page}
\markboth{Synthèse}{Synthèse}

\begin{center}
\begin{tcolorbox}[colback=typhoonDeep,colframe=typhoonBlue!0,arc=3pt,
  left=10pt,right=10pt,top=8pt,bottom=8pt,width=\textwidth]
\centering
{\color{white}\sffamily\bfseries\large \typhoon{} — données climatiques par
bâtiment}\\[3pt]
{\color{white!85}\sffamily\small Une adresse, treize aléas, un bâtiment, une
provenance sur chaque fait.}
\end{tcolorbox}
\end{center}

\section*{Le problème}

Géorisques publie les périmètres d'aléas mais ne connaît pas le bâtiment. La
BDNB décrit les bâtiments mais ne sait pas lesquels sont inondables. Aucune
source ne relie exposition et vulnérabilité à l'échelle du bâtiment. Un
souscripteur comble ce vide à l'œil, sans traçabilité.

\section*{La réponse}

Un service qui prend une adresse française et renvoie un contrat JSON
canonique : treize aléas normalisés, le passage verbatim de la BDNB, une
provenance et un niveau de résolution sur chaque fait — et aucun score.

\section*{Les trois décisions structurantes}

\begin{enumerate}[leftmargin=1.5em,itemsep=3pt]
  \item \textbf{Refuser le score composite.} Sous Solvabilité~II, un score
        tiers engage son producteur et se fait reconstruire ou refuser par
        l'actuaire. Le champ de résolution devient le signal de qualité.
  \item \textbf{Faire de l'incertitude une fonctionnalité.} Trois états
        (présent / absent / inconnu) et une distinction stricte entre absence
        de risque et panne de service.
  \item \textbf{Soustraire plutôt qu'ajouter.} La phase la plus productive du
        stage a consisté à supprimer quatre connecteurs, une couche de
        notation et un tableau de bord.
\end{enumerate}

\section*{Chiffres clés}

\begin{center}
\begin{tabular}{@{}r@{\hspace{6pt}}l@{\hspace{18pt}}r@{\hspace{6pt}}l@{}}
\textbf{357} & tests automatisés hors-ligne & \textbf{28} & exigences fonctionnelles \\
\textbf{214} & tests backend (pytest) & \textbf{18} & tâches planifiées (+10 scénarios) \\
\textbf{143} & tests frontend (Vitest) & \textbf{13} & aléas normalisés \\
\textbf{70} & fichiers Python (backend) & \textbf{4} & aléas certifiés \\
\textbf{68} & modules TypeScript (frontend) & \textbf{9} & sources externes intégrées \\
\end{tabular}
\end{center}

\section*{Résultat}

Un service déployé et vérifié en ligne le 15 septembre 2026, couvrant
l'ensemble du parcours : géocodage, jointure géométrique, fiche de
vulnérabilité, mise en situation cartographique et export. Le prototype
démontre une thèse : la donnée publique française permet de produire une
information de risque climatique par bâtiment, traçable et défendable — à
condition d'accepter de ne pas tout savoir et de le dire.

\vfill

\begin{tcolorbox}[colback=typhoonAmber!8,colframe=typhoonAmber,arc=2pt,
  left=8pt,right=8pt,top=6pt,bottom=6pt]
{\sffamily\bfseries\small Ce que le prototype n'est pas :}
{\small un produit commercialisable. L'authentification est simulée, le
magasin de lots est en mémoire, la certification attend sa qualification
humaine et le budget de latence n'est pas calibré. Ces limites sont détaillées
au chapitre~9 et assumées.}
\end{tcolorbox}
